\newpage
\newpage
  
\appendix

\section{Outlook}\label{appendix: outlook}

\subsection{Self-Improvement Mechanisms}

One limitation of the MetaGPT version in the main text of this paper is that each software project is executed independently. However, through active teamwork, a software development team should learn from the experience gained by developing each project, thus becoming more compatible and successful over time.  

This is somewhat related to the idea of recursive self-improvement, first informally proposed in 1965 \citep{DBLP:journals/ac/Good65}, with first concrete implementations since 1987 \citep{schmidhuber1987evolutionary, schmidhuber1993self, schmidhuber1998reinforcement}, culminating in the concept of mathematically optimal self-referential self-improvers  \citep{Schmidhuber:03gm3,Schmidhuber:09gm}. Generally speaking, a system should learn from experience in the real world, and meta-learn better learning algorithms from experiences of learning, and meta-meta-learn better meta-learning algorithms from experiences of meta-learning, etc., without any limitations except those of computability and physics. 


More recent, somewhat related work leverages the reasoning ability of Large Language Models (LLMs) and recursively improves prompts of LLMs, to improve performance on certain downstream tasks \citep{fernando2023promptbreeder, zelikman2023self}, analogous to the adaptive prompt engineer of 2015 \citep{schmidhuber2015learning} where one neural network learns to generate sequence of queries or prompts for another pre-trained neural network whose answers may help the first network to learn new tasks more quickly.

In our present work, we also explore a self-referential mechanism that recursively modifies the constraint prompts of agents based on information they observe during software development.
Our initial implementation works as follows. Prior to each project, every agent in the software company reviews previous feedback and makes necessary adjustments to their constraint prompts. This enables them to continuously learn from past project experiences and enhance the overall multi-agent system by improving each individual in the company. We first establish a \textit{handover feedback} action for each agent. This action is responsible for critically summarizing the information received during the development of previous projects and integrating this information in an updated constraint prompt. The summarized information is stored in \textit{long-term memory} such that it can be inherited by future constraint prompt updates. 
When initiating a new project, each agent starts with a \textit{react} action. Each agent evaluates the received feedback and summarizes how they can improve in a constraint prompt.


One current limitation is that these summary-based optimizations only modify constraints in the specialization of roles  (\Cref{sec:agents_in_sop}) rather than structured communication interfaces in communication protocols (\Cref{sec:subscription}). Future advancements are yet to be explored.

\subsection{Multi-Agent Economies}
In real-world teamwork, the interaction processes are often not hardcoded. 
For example, in a software company, the collaboration SOP may change dynamically.

One implementation of such self-organization is discussed in the paper on a ``Natural Language-Based Society of Mind'' (NLSOM)~\citep{zhuge2023mindstorms}, which introduced the idea of an ``Economy of Minds'' (EOM), a Reinforcement Learning (RL) framework for societies of LLMs and other agents. Instead of using standard RL techniques to optimize the total reward of the system through modifications of neural network parameters, EOMs use the principles of supply and demand in free markets to assign credit (money) to those agents that contribute to economic success (reward). 

%
The recent agent-based platform of DeepWisdom (AgentStore\footnote{\href{http://beta.deepwisdom.ai}{http://beta.deepwisdom.ai}}) is compatible with the credit assignment concept of EOMs. Each agent in AgentStore provides a list of services with corresponding costs. A convenient API is provided so that human users or agents in the platform can easily purchase services from other agents to accomplish their services. Figure \ref{fig:agentstore} displays the User Interface (UI) of AgentStore, where various agents with different skills are showcased. 
Besides, individual developers can participate in building new agents and enable collaborative development within the community. Specifically, AgentStore allows users to subscribe to agents according to their demands and pay according to their usage. Moreover, users can purchase additional capabilities to expand the plug-and-play functions of their existing agents. This allows users to gradually upgrade their agents. Within the MetaGPT framework, AgentStore can support the collaboration of various agents. Users can collect several agents together to carry out more complex tasks or projects, and all the agents share and comply with development and communication protocols defined in MetaGPT.  


\begin{figure}[!ht]
    \centering
    \includegraphics[width=1\textwidth]{imgs/12-agentstore.jpeg}
    \caption{AgentStore is a platform dedicated to serving users in the creation and development of agents within the MetaGPT framework. This platform provides users with an operational interface, allowing users to easily manage a variety of agents with different emotions, personalities, and capabilities for specific tasks.}
    \label{fig:agentstore}
\end{figure}


\newpage
\section{A Demo of the Execution}
\label{appendix:example}
In this section, we outline the complete process of software development using MetaGPT. 
%
It begins with a user's input command (as shown in~\Cref{listing:user_input}) and ends with software designed according to the user's specifications. 

\subsection{User Input}
Upon receiving an instruction from the user, MetaGPT collaborates with a professional development team to fulfill the task. Here is a demo of user input:
\begin{tcolorbox}[title={\textbf{\small User input}}, colback=whitesmoke, colframe=royalblue(web), boxrule=2pt, arc=0mm]
{\scriptsize
\begin{lstlisting}[style=mystyle]
write a python3 GUI app such that you can draw an image with it
\end{lstlisting}
}
\end{tcolorbox}
\label{listing:user_input}

\subsection{MetaGPT development process}
Now we provide a step-by-step explanation of the standardized output process for each agent.

% \newpage
\paragraph{Product Manager}
The Product Manager generates a Product Requirement Document (PRD), as detailed in the specified documentation. This document encompasses goals, user stories, competitive analysis, requirement analysis and requirement
pool. Additionally, a competitive quadrant chart is produced (see~\Cref{fig:quadrantChart}). Subsequently, these documents and charts are handed over to the architect for system design.


% \begin{tcolorbox}[title={\textbf{\small Product Requirement Document}}, colback=white, colframe=royalblue(web), boxrule=2pt, arc=0mm]
% {\scriptsize
% \begin{lstlisting}[style=PythonCode]
\begin{tcolorbox}[title={\textbf{\small Product Requirement Document}}, colback=whitesmoke, colframe=royalblue(web), boxrule=2pt, arc=0mm]
{\scriptsize
\begin{lstlisting}[style=mystyle]
## Original Requirements
The boss requires a Python3 GUI color meter that outputs the RGB values when it moves.

## Product Goals
```python
[
    "Create a user-friendly GUI color meter",
    "Ensure the color meter accurately outputs RGB values",
    "Ensure the color meter updates RGB values in real-time as it moves"
]
```

## User Stories
```python
[
    "As a user, I want to select any color on the screen, so that I can get its RGB values",
    "As a user, I want the RGB values to update in real-time as I move the color meter, so that I can get accurate color information",
    "As a user, I want a simple and intuitive interface, so that I can easily use the color meter"
]
```

## Competitive Analysis
```python
[
    "Color Cop: A popular color picker for Windows. It's simple and easy to use, but lacks real-time RGB value updates",
    "Just Color Picker: Offers real-time color picking and multiple color code formats. However, its interface is cluttered and not very user-friendly",
    "Instant Eyedropper: A simple color picker that lives in the system tray. It's convenient but lacks advanced features",
    "ColorZilla: A color picker extension for browsers. It's convenient for web design but not suitable for desktop applications",
    "ColorPic: Offers color palettes and mixer tools. It's feature-rich but can be overwhelming for simple tasks"
]
```

\end{lstlisting}
}
\end{tcolorbox}

\begin{tcolorbox}[title={\textbf{\small Product Requirement Document}}, colback=whitesmoke, colframe=royalblue(web), boxrule=2pt, arc=0mm]
{\scriptsize
\begin{lstlisting}[style=mystyle]

## Requirement Analysis
The product should be a Python3 GUI application that serves as a color meter. It should allow users to select any color on the screen and display the corresponding RGB values. The RGB values should update in real-time as the color meter moves. The interface should be simple and user-friendly.

## Requirement Pool
```python
[
    ("Design a user-friendly GUI for the color meter", "P0"),
    ("Implement color selection functionality", "P0"),
    ("Display RGB values of the selected color", "P0"),
    ("Update RGB values in real-time as the color meter moves", "P0"),
    ("Test the application for accuracy and performance", "P1")
]
```

## UI Design draft
The UI should be minimalistic and intuitive. It should primarily consist of a color meter, which could be a magnified area of the screen that the user can move around. The RGB values of the currently selected color should be prominently displayed, updating in real-time as the color meter moves. The layout should be clean and uncluttered, with a focus on functionality and ease of use.

## Anything UNCLEAR
There are no unclear points.
\end{lstlisting}
}
\end{tcolorbox}





\label{supp:listing:product_manager}

\begin{figure}[!ht]
    \centering
    \includegraphics[width=0.6\textwidth]{imgs/6-quadrant_chart.png}
    \caption{The quadrant chart for the “Drawing App” generated by the \emph{product manager} agent.}
    \label{fig:quadrantChart}
\end{figure}

\newpage

\paragraph{Architect}
Based on the requirements in PRD, the Architect agent devises technical specifications including system architecture diagrams and interface definitions. Initially, the Architect defines the overarching technical trajectory. Subsequently, the project's architecture, including files, classes (\Cref{fig:data_structures}) and the sequence flow chart (\Cref{fig:program_call_flow}), is designed. The Architect's documentation is then given to the project manager for task allocation and execution.
\begin{tcolorbox}[title={\textbf{\small System design}}, colback=whitesmoke, 
colframe=royalblue(web), boxrule=2pt, arc=0mm]
{\scriptsize
\begin{lstlisting}[style=mystyle]
## Implementation approach
We will use Python's Tkinter library to create the GUI for the color meter. Tkinter is a standard Python interface to the Tk GUI toolkit and is included with most Python installations. It is simple to use and powerful enough for our needs.

The color selection functionality will be implemented using the PIL (Pillow) library, which allows us to capture the screen's pixel data. By moving the mouse, we can get the RGB values of the pixel under the mouse cursor in real-time.

The RGB values will be displayed on the GUI using Tkinter's Label widget. We will use Python's threading library to update the RGB values in real-time as the color meter moves.

The application will be tested for accuracy and performance using Python's unittest library.

## Python package name
```python
"color_meter"
```

## File list
```python
[
    "main.py",
    "color_picker.py",
    "gui.py",
    "tests.py"
]
```

## Anything UNCLEAR
The requirement is clear to me.
\end{lstlisting}
}
\end{tcolorbox}

\vspace{30pt}
\begin{figure}[!ht]
    \centering
    \includegraphics[width=0.47\textwidth]{imgs/7-data_structures_and_interface_definitions.png}
    \caption{Data structures and interface definitions for the “Drawing App” generated by the \emph{architect} agent.}
    \vspace{30pt}
    \label{fig:data_structures}
\end{figure}


\paragraph{Project Manager}
The Project Manager breaks down the project into a task list. Furthermore, each code file is analyzed based on its intended functionality and then treated as a separate task assigned to Engineers.

\begin{figure}[!ht]
    \centering
    \includegraphics[width=0.83\textwidth]{imgs/8-program_call_flow.png}
    \caption{The program call flow for the “Drawing App” generated by the \emph{architect} agent.}
    \label{fig:program_call_flow}
    %\vspace{50pt}
\end{figure}

\newpage


\begin{tcolorbox}[title={\textbf{\small Tasks}}, colback=whitesmoke, colframe=royalblue(web), boxrule=2pt, arc=0mm]
{\scriptsize
\begin{lstlisting}[style=mystyle]
## Required Python third-party packages

```python
"""
tkinter==8.6
Pillow==8.3.1
"""
```

## Required Other language third-party packages
```python
"""
No third-party packages required in other languages.
"""
```

## Full API spec
```python
"""
No APIs are used in this project.
"""
```

## Logic Analysis
```python
[
    ("main.py", "Contains the main application loop and initializes all the other components."),
    ("canvas.py", "Implements the drawing canvas, responds to mouse events, and uses the selected tool and color."),
    ("tools.py", "Implements the tool selection interface."),
    ("color_picker.py", "Implements the color selection interface."),
    ("file_manager.py", "Implements the file saving and opening functionality.")
]
```
\end{lstlisting}
}
\end{tcolorbox}

\begin{tcolorbox}[title={\textbf{\small Product Requirement Document}}, colback=whitesmoke, colframe=royalblue(web), boxrule=2pt, arc=0mm]
{\scriptsize
\begin{lstlisting}[style=mystyle]
## Task list
```python
[
    "main.py",
    "canvas.py",
    "tools.py",
    "color_picker.py",
    "file_manager.py"
]
```

## Shared Knowledge
```python
"""
The 'Tkinter' library is used for creating the GUI of the application. It provides various widgets like frames and canvases which are used in this project.
The 'Pillow' library is used for handling image files. It is used in the 'file_manager.py' for saving and opening image files.
The 'main.py' file is the entry point of the application. It initializes all the other components and starts the application loop.
"""
```

## Anything UNCLEAR
The project requirements and implementation approach are clear. However, we need to ensure that all team members are familiar with the Tkinter and Pillow libraries. If not, they may need some time to learn these libraries before starting the development.
\end{lstlisting}
}
\end{tcolorbox}



\newpage
\paragraph{Engineer}
Given the provided file structure and function definitions, an Engineer agent requires only fundamental development skills to complete the development tasks. Due to the large number of files, we present only one auto-generated code file here.% as an illustration.

\begin{tcolorbox}[title={\textbf{\small Generated code (color\_picker.py)}}, 
colback=whitesmoke, colframe=royalblue(web), boxrule=2pt, arc=0mm]
{\scriptsize
\begin{lstlisting}[style=mystyle]
import tkinter as tk
from tkinter import colorchooser
class ColorPicker:
    def __init__(self, root: tk.Tk, color: str = 'black'):
        self.frame = tk.Frame(root)
        self.color = color

        self.color_button = tk.Button(self.frame, text='Color', command=self.select_color)
        self.color_button.pack(fill='x')

    def select_color(self):
        color = colorchooser.askcolor()[1]
        if color is not None:
            self.color = color

    def pack(self, **kwargs):
        self.frame.pack(**kwargs)

    def get_color(self):
        return self.color
\end{lstlisting}
}
\end{tcolorbox}


\paragraph{QA Engineer}
Upon receiving the code output from the Engineer, the QA Engineer generates unit test code and reviews it to identify and fix any bugs, ensuring high-quality software.
\begin{tcolorbox}
[title={\textbf{\small Generated Unit tests (test.py)}}, 
colback=whitesmoke, colframe=royalblue(web), boxrule=2pt, arc=0mm]
{\scriptsize
\begin{lstlisting}[style=mystyle]
import unittest
import tkinter as tk
# You might want to move the ColorPicker class to another file, let's say color_picker.py
# from color_picker import ColorPicker

class TestColorPicker(unittest.TestCase):

    def setUp(self):
        self.root = tk.Tk()
        self.cp = ColorPicker(self.root)

    def tearDown(self):
        # This is important to close the tk window after tests
        self.root.destroy()

    def test_initial_color(self):
        # Test the default color during initialization
        self.assertEqual(self.cp.get_color(), 'black')

    def test_set_and_get_color(self):
        # Test setting a new color and retrieving it
        new_color = '#ffffff'  # white color
        self.cp.color = new_color
        self.assertEqual(self.cp.get_color(), new_color)


if __name__ == '__main__':
    unittest.main()
\end{lstlisting}
}
\end{tcolorbox}




\paragraph{Output}
Ultimately, as shown in~\Cref{fig:demo}, MetaGPT generates a functional application named ``Drawing App".
\begin{figure}[!ht]
    \centering
    \includegraphics[width=\textwidth]{imgs/9-drawing_app_demo.png}
    \caption{The ``Drawing App" generated by MetaGPT.}
    \label{fig:demo}
\end{figure}

\newpage
\section{Experiments}\label{appendix: exp}

\subsection{Details of the SoftwareDev Dataset}
The SoftwareDev dataset includes 70 diverse software development tasks. \Cref{tab:task_details} displays the names and detailed prompts of 11 tasks within the dataset. Note that the first seven tasks listed are used in the main experiments of this paper. 

\subsection{Additional results}
\paragraph{Quantitative results of MetaGPT}
% 
As shown in~\Cref{tab:task_executability}, MetaGPT achieves an average score of 3.9, surpassing ChatDev's score of 2.1 \cite{zhao2023chat}, which is based on the Chat chain. Compare the scores of general intelligent algorithms, including AutoGPT~\cite{autogpt_github}, which all score 1.0, failing to generate executable code. 
We observe that the generated code is often short, lacks comprehensive logic, and tends to fail to handle cross-file dependencies correctly.

While models such as AutoGPT~\citep{autogpt_github}, Langchain~\citep{Chase_LangChain_2022}, and AgentVerse~\citep{agentverse_github} display robust general problem-solving capabilities, they lack an essential element for developing complex systems: systematically deconstructing requirements.
% 
Conversely, MetaGPT simplifies the process of transforming abstract requirements into detailed class and function designs through a specialized division of labor and SOPs workflow.
%
When compared to ChatDev~\citep{zhao2023chat}, MetaGPT's structured messaging and feedback mechanisms not only reduce loss of communication information but also improve the execution of code.
% 
\begin{table}[htb]
  \centering
      \caption{\textbf{Executability comparison.} The executability scores are on a grading system ranging from '1' to '4'. A score of '1' signifies complete failure, '2' denotes executable code, '3' represents largely satisfying expected workflow, and '4' indicates a perfect match with expectations.}
      \label{tab:task_executability}
\renewcommand\tabcolsep{8pt}
\renewcommand\arraystretch{1.2}
  \begin{tabular}{l|ccccc} %
\Xhline{1.5pt}
\rowcolor{paleaqua} 
\textbf{Task} & \textbf{AutoGPT} & \textbf{LangChain} & \textbf{AgentVerse}& \textbf{ChatDev} & \textbf{MetaGPT} \\
\Xhline{1.5pt}
Flappy bird&1&1&1&2&\textbf{3} \\
\rowcolor{gray!20} Tank battle game&1&1&1&2&\textbf{4} \\
2048 game&1&1&1&1&\textbf{4} \\
\rowcolor{gray!20} Snake game&1&1&1&3&\textbf{4} \\
Brick breaker game&1&1&1&1&\textbf{4} \\
\rowcolor{gray!20} Excel data process&1&1&1&\textbf{4}&\textbf{4} \\
CRUD manage&1&1&1&2&\textbf{4} \\
\Xhline{1pt}
\rowcolor{gray!20} Average score &1.0&1.0 &1.0&2.1&\textbf{3.9} \\
\Xhline{1.5pt}
\end{tabular}
\end{table}
\paragraph{Quantitative results of MetaGPT w/o executable feedback}
\Cref{tab:result_details} presents the performance of MetaGPT with GPT-4 32K on 11 tasks within the SoftwareDev dataset. It also shows the average performance across all 70 tasks (in the last line). Note that the version of MetaGPT used here is the basic version without the executable feedback mechanism.

\paragraph{Quantitative results of MetaGPT with different LLMs}
To verify the performance of MetaGPT on different LLMs, we randomly selected 5 SoftwareDev tasks and conducted experiments using GPT-3.5 and Deepseek Coder 33B\footnote{\href{https://deepseekcoder.github.io}{https://deepseekcoder.github.io}} as backends. As shown in \Cref{tab:MetaGPT_with_LLMs}, the results indicate that although MetaGPT can complete tasks with these LLMs, using GPT-4 as the backend yields superior performance.

\begin{table}[htb]
  \centering
      \caption{\textbf{Performance of MetaGPT on SoftwareDev using different LLMs as agent backends.} }
      \label{tab:MetaGPT_with_LLMs}
\renewcommand\tabcolsep{2.5pt}
\renewcommand\arraystretch{1.2}
% \resizebox{\textwidth}{!}{%
  \begin{tabular}{c|ccccc} %
\Xhline{1.5pt}
\rowcolor{paleaqua} 
\textbf{Model} & \textbf{Open source} & \textbf{Time(/s)} & \textbf{\# Lines}& \textbf{Executability} & \textbf{Revisions} \\
\Xhline{1.2pt}
MetaGPT (w/ GPT-3.5)& \XSolidBrush &75.18&161.6&2.8&\ 2.4 \\
\rowcolor{gray!20} MetaGPT (w/ GPT-4)& \XSolidBrush &552.94&178.2&3.8&\ 1.2 \\
MetaGPT (w/ Deepseek Coder 33B)& \CheckmarkBold &1186.20&120.2&1.4&\ 2.6 \\
\Xhline{1.2pt}
\end{tabular}
% }
\end{table}


\paragraph{Impact of Instruction Levels (High-level \textit{v.s.} Detailed Instructions)}
Does the variation in the level of initial input from humans significantly influence performance outcomes?  For examples:
\begin{enumerate}
    \item \textbf{High-level prompt}: Create a brick breaker game.
    \item \textbf{Detailed prompt}: Creating a brick breaker game. In a brick breaker game, the player typically controls a paddle at the bottom of the screen to bounce a ball towards a wall of bricks. The goal is to break all the bricks by hitting them with the ball.
\end{enumerate}

Additional experiments were conducted to investigate this aspect: we selected 5 tasks from SoftwareDev, and constructed detailed prompts for them. Here are the experimental results:
\begin{table}[htb]
  \centering
      \caption{\textbf{Impact of Instruction Levels.} The executability is scored on a grading system ranging from `1' to `4'. A score of `1' signifies complete failure, `2' denotes runnable code, `3' represents largely expected workflow, and `4' indicates a perfect match to expectations.}
      \label{tab:impact_of_instruction_levels}
\renewcommand\tabcolsep{3pt}
\renewcommand\arraystretch{1.2}
% \resizebox{\columnwidth}{!}{%
\begin{tabular}{l|ccccccc} %
\Xhline{1.5pt}
\rowcolor{paleaqua} 
\textbf{Model} & \textbf{\# Word} & \textbf{Time(/s)} & \textbf{Token usage}& \textbf{\# Lines} & \textbf{Executability} & \textbf{Productivity} & \textbf{Reversions} \\
\Xhline{1.5pt}
High-level&13.2&552.9&28384.2&178.2&3.8 & 163.8 & 1.2 \\
\rowcolor{gray!20} Detailed&42.2&567.8&29657.0&257.0&4.0&118.0 & 1.6 \\
\Xhline{1.5pt}
\end{tabular}
% }
\end{table}

We observe that: detailed prompts lead to better software projects with lower productivity ratios because of clearer requirements and functions, while simple inputs can still generate good enough software using MetaGPT with an executability rating of 3.8, which is comparable to the detailed prompt scenario. (Note that, Productivity = Token usage / Total Code Lines. The lower this ratio, the better.)

\paragraph{The performance of GPT variants in HumanEval benchmark}
We use the GPT-4's 67\% HumanEval score~\citep{openai2023gpt4} as our baseline, acknowledging its acceptance in the HumanEval benchmark. We further extend to experiments(five times) with GPT-4 (gpt-4-0613) and GPT-3.5-Turbo (gpt-3.5-turbo-0613) under various conditions to assess performance. \textbf{(A)}~We directly called the OpenAI API with the prompt in HumanEval. \textbf{(B)}~We called the OpenAI API and parsed the code with regex in the response. \textbf{(C)}~We added an additional system prompt, then called the OpenAI API. The prompt is "You are an AI that only responds with Python code, NOT ENGLISH. You will be given a function signature and its docstring by the user. Write your full implementation (restate the function signature)." As 
shown in ~\Cref{tab:performance_on_gpt}, GPT-4 is more sensitive to prompt, code parser, and post-processing results on the HumanEval data set. It is difficult for GPT-3.5-Turbo to return the correct completion code without prompt words.


\begin{table}[htb]
  \centering
      \caption{\textbf{Performance of GPT models on HumanEval.}{
      Experiments were conducted five times using gpt-4-0613 and gpt-3.5-turbo-0613 with different settings.
      % Comparative Analysis of GPT-4 and GPT-3.5-Turbo Performance on HumanEval Benchmark Across Different Experimental Conditions.
      }}
      \label{tab:performance_on_gpt}
\renewcommand\tabcolsep{7.2pt}
\renewcommand\arraystretch{1.2}
% \resizebox{\columnwidth}{!}{%
  \begin{tabular}{l|cccccccc} %
\Xhline{1.5pt}
\rowcolor{paleaqua} 
\textbf{Settings} & \textbf{Model} & \textbf{1} & \textbf{2}& \textbf{3} & \textbf{4} & \textbf{5} & \textbf{Avg.} & \textbf{Std.} \\
\Xhline{1.5pt}
A&gpt-4-0613&0.732&0.707&0.732 &0.713 & 0.738 & 0.724 & 0.013\\
\rowcolor{gray!20} A &gpt-3.5-turbo-0613&0.360&0.366&0.360& 0.348 & 0.354 & 0.357 & 0.007\\
B&gpt-4-0613&0.787&0.811&0.817&0.829& 0.817 & 0.812 & 0.016 \\
\rowcolor{gray!20} B &gpt-3.5-turbo-0613&0.348&0.354&0.348&0.335 & 0.348& 0.346& 0.007\\
C&gpt-4-0613&0.805&0.805&0.817&0.793&0.780 &0.800 & 0.014 \\
\rowcolor{gray!20} C &gpt-3.5-turbo-0613&0.585&0.567&0.573&0.579&0.579 & 0.577& 0.007 \\
\Xhline{1.5pt}
\end{tabular}
% }

\end{table}


\paragraph{Qualitative results}
\Cref{fig:system_interface_design_of_recommendation_engine} and \Cref{fig:program_call_flow} illustrate the outcomes of the Architect agent's efforts to design a complex recommender system.
%
These figures showcase the comprehensive system interface design and program call flow. 
%
The latter is essential for creating a sophisticated automated system. It is crucial to emphasize the importance of this division of labor in developing an automated software framework.

\section{Limitation and Ethics Concerns}

\subsection{Limitation}
\paragraph{System side} At present, our system cannot fully cater to specific scenarios, such as UI and front-end, as we have yet to incorporate such agents and multimodal tools. Furthermore, despite generating the most amount of code among comparable frameworks, it remains challenging to fulfill real-world applications' diverse and complex requirements.

\paragraph{Human user side} A key challenge for users is to interrupt the running process of each agent, or set the starting running point (checkpoint) for each agent.

\subsection{Ethics Concerns}
\paragraph{Unemployment and Skill Obsolescence} MetaGPT enables more people to program in natural languages, thereby making it easier for engineers to get started. Over the years, programming languages have evolved from punched cards to assembly, C, Java, Python, and now natural language. As a result, humans have become more proficient at programming, increasing the demand for programming-related positions. Furthermore, programming with natural language may offer a significantly easier learning curve, making programming more accessible to a broader audience.

\paragraph{Transparency and Accountability} MetaGPT is an open-source framework that facilitates interactive communication between multiple agents through natural language. Humans can initiate, observe, and stop running with the highest level of control. It provides real-time interpretation and operation of the natural language, displayed on the screen and logs, ensuring transparency. MetaGPT enhances ``natural language programming" capabilities, and human engineers are the users and responsible for the outcomes.

\paragraph{Privacy and Data Security} MetaGPT operates locally, ensuring user data privacy and security. It does not collect user data. For interactions with third-party LLMs, such as those by OpenAI, users are encouraged to refer to the respective privacy policies (e.g., OpenAI Privacy Policy). However, we provide the option of open-source LLMs as backends.

\begin{figure}[t]
    \centering
     \includegraphics[width=\textwidth]{imgs/10-system_interface_design_of_recommendation_engine.png}
     \vspace{-15pt}
    \caption{
    The system interface design for ``recommendation engine development'' is generated by the \emph{architect} agent (\textbf{zoom in for a better view}).
    }
    \label{fig:system_interface_design_of_recommendation_engine}
\end{figure}

\begin{figure}[t]
    \centering
     \includegraphics[width=\textwidth]{imgs/11-program_call_flow_for_recommendation_engine_development.png}
     \vspace{-15pt}
    \caption{
    The program call flow for ``recommendation engine development'' generated by the \emph{architect} agent (\textbf{zoom in for a better view}).
    }
    \label{fig:program_call_flow}
\end{figure}

\begin{table}[htb]
\centering
\caption{Examples of SoftwareDev dataset.}
\label{tab:task_details}
\renewcommand\tabcolsep{5.2pt}
\renewcommand\arraystretch{1.2}
\resizebox{\textwidth}{!}{%
\LARGE
\begin{tabular}{llp{18.5cm}}
\hline

\hline

\hline

\hline
Task ID & Task & Prompt \\

\hline

\hline
0 & Snake game & Create a snake game. \\ % \midrule
1 & Brick breaker game & Create a brick breaker game. \\ % \midrule
2 & 2048 game & Create a 2048 game for the web. \\ % \midrule
3 & Flappy bird game & Write p5.js code for Flappy Bird where you control a yellow bird continuously flying between a series of green pipes. The bird flaps every time you left click the mouse. If it falls to the ground or hits a pipe, you lose. This game goes on indefinitely until you lose; you get points the further you go. \\ % \midrule
4 & Tank battle game & Create a tank battle game. \\ % \midrule
5 & Excel data process & Write an excel data processing program based on streamlit and pandas. The screen first shows an excel file upload  button. After the excel file is uploaded, use pandas to display its data content. The program is required to be concise, easy to maintain, and not over-designed. It uses streamlit to process web screen displays, and pandas is sufficient to process excel reading and display. Please make sure others can execute directly without introducing additional packages.
%Others do not need to introduce additional packages. 
%juergen what do you mean by: Others do not need to introduce additional packages.???
\\ % \midrule
6 & CRUD manage & Write a management program based on the crud addition, deletion, modification and query processing of the customer business entity. The customer needs to save this information: name, birthday, age, sex, and phone. The data is stored in client.db, and there is a judgement whether the customer table exists. If it doesn't, it needs to be created first. Querying is done by name; same for deleting. The program is required to be concise, easy to maintain, and not over-designed. The screen is realized through streamlit and sqlite---no need to introduce other additional packages. \\ % \midrule
\hline

\hline
7 & Music transcriber & Develop a program to transcribe sheet music into a digital format; providing error-free transcribed symbolized sheet music intelligence from audio through signal processing involving pitch and time slicing then training a neural net to run Onset Detected CWT transforming scalograms to chromagrams decoded with Recursive Neural Network focused networks. \\ % \midrule
8 & Custom press releases & Create custom press releases; develop a Python script that extracts relevant information about company news from external sources, such as social media; extract update interval database for recent changes. The program should create press releases with customizable options and export writings to PDFs, NYTimes API JSONs, media format styled with interlink internal fixed character-length metadata. \\ % \midrule
9 & Gomoku game & Implement a Gomoku game using Python, incorporating an AI opponent with varying difficulty levels. \\ % \midrule
10 & Weather dashboard & Create a Python program to develop an interactive weather dashboard. \\ 
\hline

\hline

\hline

\hline

\hline
\end{tabular}
}
\end{table}


\section{More Discussions}
\subsection{Deep-seated challenges}
MetaGPT also alleviates or solves these challenges with its unique designs:
\paragraph{Use Context Efficiently} Two sub-challenges are present. First, unfolding short natural language descriptions accurately to eliminate ambiguity. Second, maintaining information validity in lengthy contexts, enables LLMs to concentrate on relevant data without distraction.
\paragraph{Reduce Hallucinations} Using LLMs to generate entire software programs faces code hallucination problems—-including incomplete implementation of functions, missing dependencies, and potential undiscovered bugs, which may be more serious. %Deep-seated challenges in task decomposition and translating from natural language to code language persist.
LLMs often struggle with software generation due to vague task definitions. Focusing on granular tasks like requirement analysis and package selection offers guided thinking, which LLMs lack in broad task solving.


\subsection{Information Overload}
In MetaGPT, we use a global message pool and a subscription mechanism to address ``information overload," which refers to the problem of receiving excessive or irrelevant information. This issue is dependent on specific applications. MetaGPT employs a message pool to streamline communication, ensuring efficiency. Additionally, a subscription mechanism filters out irrelevant contexts, enhancing the relevance and utility of the information. This design is particularly crucial in software design scenarios and standard operating procedures (SOPs) where effective communication is essential.

\clearpage

\begin{sidewaystable}[htb]
\centering
\caption{\textbf{Additional results of pure MetaGPT w/o feedback on SoftwareDev.} Averages (Avg.) of 70 tasks are calculated and 10 randomly selected tasks are included. `\#' denotes `The number of', while `ID' is `Task ID'.}
\label{tab:result_details}
\renewcommand\tabcolsep{0.8pt}
\renewcommand\arraystretch{1.2}
% \resizebox{\textwidth}{!}{%
\small
\begin{tabular}{lccccccccccp{2.5cm}c}
\hline

\hline

\hline

\hline
ID & \multicolumn{3}{c}{Code statistics} & \multicolumn{3}{c}{Doc statistics} & & \multicolumn{3}{c}{Cost statistics} & Cost of revision & Code executability
\\
\cline{2-3} \cline{4-7} \cline{8-11} 
& \#code files & \#lines of code & \#lines per code file & \#doc files & \#lines of doc & \#lines per doc file & \#prompt tokens & \#completion tokens & time costs & money costs & &  \\ % \midrule
\hline

\hline
0 & 5.00 & 196.00 & 39.20 & 3.00 & 210.00 & 70.00 & 24087.00 & 6157.00 & 582.04 & \$~1.09 & 1. TypeError & 4 \\ % \midrule
1 & 6.00 & 191.00 & 31.83 & 3.00 & 230.00 & 76.67 & 32517.00 & 6238.00 & 566.30 & \$~1.35 & 1. TypeError & 4 \\ % \midrule
2 & 3.00 & 198.00 & 66.00 & 3.00 & 235.00 & 78.33 & 21934.00 & 6316.00 & 553.11 & \$~1.04 & 1. lack @app.route('/') & 3 \\ % \midrule
3 & 5.00 & 164 & 32.80 & 3.00 & 202.00 & 67.33 & 22951.00 & 5312.00 & 481.34 & \$~1.01 & 1. PNG file missing 2. Compile bug fixes & 2 \\ % \midrule
4 & 6.00 & 203.00 & 33.83 & 3.00 & 210.00 & 70.00 & 30087.00 & 6567.00 & 599.58 & \$~1.30 & 1. PNG file missing 2. Compile bug fixes 3. pygame.surface not initialize & 3 \\ % \midrule
5 & 6.00 & 219.00 & 36.50 & 3.00 & 294.00 & 96.00 & 35590.00 & 7336.00 & 585.10 & \$~1.51 & 1. dependency error 2. ModuleNotFoundError & 4 \\ % \midrule
6 & 4.00 & 73.00 & 18.25 & 3.00 & 261.00 & 87.00 & 25673.00 & 5832.00 & 398.83 & \$~0.90 & 0 & 4 \\ % \midrule
\hline

\hline
7 & 4.00 & 316.00 & 79.00 & 3.00 & 332.00 & 110.67 & 29139.00 & 7104.00 & 435.83 & \$~0.92 & 0 & 4 \\ % \midrule
8 & 5.00 & 215.00 & 43.00 & 3.00 & 301.00 & 100.33 & 29372.00 & 6499.00 & 621.73 & \$~1.27 & 1. tensorflow version error 2. model training method not implement & 2 \\ % \midrule
9 & 5.00 & 215.00 & 43.00 & 3.00 & 270.00 & 90.00 & 24799.00 & 5734.00 & 550.88 & \$~1.27 & 1. dependency error 2. URL 403 error & 3 \\ % \midrule
10 & 3.00 & 93.00 & 31.00 & 3.00 & 254.00 & 84.67 & 24109.00 & 5363.00 & 438.50 & \$~0.92 & 1. dependency error 2. missing main func. & 4 \\ % \midrule
\hline

\hline
\rowcolor[gray]{.9} 
Avg. & 4.71 & 191.57 & 42.98 & 3.00 & 240.00 & 80.00 & 26626.86 & 6218.00 & 516.71 & \$1.12 & 0.51~(only consider item scored 2, 3 or 4) & 3.36 \\ 
\hline

\hline

\hline

\hline
\end{tabular}
% }
\end{sidewaystable}